\section{Respuesta a incidentes y revisión post mortem}
\subsection{Proceso de atribución de fallas}
Cada vez que un Agente falla, el equipo inicia de inmediato una revisión post mortem: primero extrae la `failure matrix`, que incluye el comando, el archivo que la desencadenó, el entorno de ejecución y el código de retorno; después compara la estrategia actual con las estrategias históricas para determinar si la regresión es del prompt, de la cadena de herramientas o del entorno de ejecución.

\begin{table}[H]
\centering
\begin{tabular}{p{0.3\textwidth} p{0.3\textwidth} p{0.3\textwidth}}
\toprule
Categoría & Condición desencadenante típica & Acción de revisión post mortem \\
\midrule
Falla del prompt & La generación no coincide con el contexto & Rediseñar el prompt como plantilla y agregar nuevos examples \\
\midrule
Error de la cadena de herramientas & El comando o la prueba se cuelga & Revisar la imagen de docker, las versiones de las dependencias y los permisos \\
\midrule
Diferencia en el entorno de evaluación & Pasa en local pero falla en CI & Comparar las variables de entorno y la versión de los datos, y hacer pruebas de dev contra prod \\
\bottomrule
\end{tabular}
\caption{Categorías de fallas y acciones de revisión post mortem.}
\end{table}

\subsection{Difusión de las minutas de revisión post mortem}
El resultado de la revisión post mortem se escribe en la wiki del equipo en formato Markdown, con la fecha, la persona responsable, la conclusión de la revisión y las acciones posteriores, y se envía un recordatorio por el canal de Slack \texttt{\#agent-postmortem}. Cada revisión post mortem también incluye capturas de pantalla clave (como el log de las pruebas o el diff) para ubicar el problema rápidamente en el futuro.

\begin{warningbox}{No dejes que la revisión post mortem se pierda}
Las minutas de la revisión post mortem deben publicarse dentro de las primeras 24 horas; si se tarda demasiado, el escenario se olvida.
\end{warningbox}

\begin{knowledgebox}{Valor de la revisión post mortem}
La revisión post mortem no solo sirve para encontrar errores: convierte las muestras de falla del Agente en material didáctico para entrenar prompts, y transforma el «por qué falló» en un patrón reproducible.
\end{knowledgebox}

\subsection{Resumen de la sección}
\begin{itemize}
    \item La atribución rápida y la revisión post mortem evitan que problemas similares se repitan.
    \item Compartir los resultados de la revisión post mortem con el equipo mejora la madurez general de la ingeniería de Agentes.
\end{itemize}

